\section{Free-choice nets}

% ── slides p.2–p.3 ──────────────────────────────────────────
For arbitrary Petri nets the analysis algorithms are expensive when they exist at all, and no general theory links structure to behaviour. The obstacle, in one phrase, is the \emph{interference between conflict and synchronisation}.\footnote{Desel J., Esparza J., \emph{Free Choice Petri Nets}, Cambridge University Press, \url{https://www7.in.tum.de/~esparza/bookfc.html}.}
\begin{center}
\begin{tikzpicture}[node distance=0.7cm and 0.9cm]
  \node[bpmn event] (p3in) {};
  \fill[accent] (p3in.center) circle (2pt);
  \node[bpmn task, rounded corners=0pt, below=of p3in, minimum width=0.8cm] (t3) {\scriptsize $t_3$};
  \node[bpmn event, below=of t3] (pmid) {};
  \node[bpmn event, left=1.2cm of pmid] (p1) {};
  \fill[accent] (p1.center) circle (2pt);
  \node[bpmn task, rounded corners=0pt, below left=0.6cm and 0.1cm of pmid, minimum width=0.8cm] (t1) {\scriptsize $t_1$};
  \node[bpmn task, rounded corners=0pt, below right=0.6cm and 0.1cm of pmid, minimum width=0.8cm] (t2) {\scriptsize $t_2$};
  \draw[bpmn flow] (p3in) -- (t3);
  \draw[bpmn flow] (t3) -- (pmid);
  \draw[bpmn flow] (p1) -- (t1);
  \draw[bpmn flow] (pmid) -- (t1);
  \draw[bpmn flow] (pmid) -- (t2);
\end{tikzpicture}
\end{center}
$t_1$ and $t_2$ are not initially in conflict ($t_2$ still waits for $t_3$'s token) but once $t_3$ fires they compete on the shared place: an uncontrollable upstream synchronisation silently induces a downstream conflict. Free-choice nets rule this out structurally.

% ── slides p.4–p.12 ──────────────────────────────────────────
\begin{definitionbox}{Free-choice nets}
\emph{Restrictive formulation}: a net is \textbf{free-choice} if for every arc $(p, t) \in F$ either $p^\bullet = \{t\}$ (no conflict on $p$) or ${}^\bullet t = \{p\}$ (no synchronisation at $t$): conflict and synchronisation never touch.
\emph{Standard formulation} (adopted here, slightly more general): for any two transitions, ${}^\bullet t_1 \cap {}^\bullet t_2 = \emptyset$ or ${}^\bullet t_1 = {}^\bullet t_2$: pre-sets are disjoint or equal (equivalently on post-sets of places). A system is free-choice if its net is. The choice between conflicting transitions is then genuinely \emph{free}: made locally on the shared pre-set, uninfluenced by the rest of the system.
\end{definitionbox}

The check is visual: whenever two transitions share an input place, do they share \emph{all} their inputs? A net where $t_1$ needs $\{p_1, p_3\}$ while $t_2$ needs only $\{p_3\}$ is not free-choice; making both consume exactly the same pre-set repairs it. Every \textbf{S-net} is free-choice (all pre-sets are singletons: equal or disjoint), and so is every \textbf{T-net} (all post-sets of places are singletons). The inclusions are strict, and the four regions of the diagram are all inhabited: a pure sequence (S and T), a free choice between two transitions (S only), an AND-join (T only), and a net combining conflicts and synchronisations kept apart (free-choice, neither S nor T):
\begin{center}
\begin{tikzpicture}
  \draw[accent, thick] (0,0) ellipse (3.0cm and 1.6cm);
  \node[accent] at (0,1.2) {\small\sffamily Free-choice};
  \draw[accent, thick] (-0.9,-0.1) ellipse (1.4cm and 0.9cm);
  \draw[accent, thick] (0.9,-0.1) ellipse (1.4cm and 0.9cm);
  \node[accent] at (-1.4,-0.1) {\small\sffamily S-nets};
  \node[accent] at (1.4,-0.1) {\small\sffamily T-nets};
\end{tikzpicture}
\end{center}

\begin{notebox}{Fundamental property; $N^*$}
\textbf{Proposition.} In a free-choice system, if a reachable $M$ enables some $t \in p^\bullet$, it enables \emph{every} $t' \in p^\bullet$ (free-choice gives ${}^\bullet t = {}^\bullet t'$, and all those places are marked). Enabledness is a property of the conflict set, not of the single transition: all or none.
\par\smallskip
\textbf{Proposition.} A workflow net $N$ is free-choice iff $N^*$ is: the reset transition's pre-set $\{o\}$ is disjoint from every other pre-set ($o$ is a sink in $N$), so adding or removing it cannot break the condition.
\end{notebox}

% ── slides p.15–p.21 ──────────────────────────────────────────
\begin{notebox}{Liveness $=$ place-liveness on free-choice systems}
A place $p$ is \textbf{live} when from every reachable marking some further reachable marking puts a token on $p$ ($\forall M \in [M_0\rangle\ \exists M' \in [M\rangle.\ M'(p) \geq 1$); the system is \textbf{place-live} when every place is live. Equivalently, $p$ is dead when it can permanently run dry, and then every transition in ${}^\bullet p \cup p^\bullet$ is dead too, since none can ever fire again through $p$. Place-liveness is the token-level analogue of liveness (which asks the same of every \emph{transition}).
\par\smallskip
In any system liveness implies place-liveness. On free-choice systems the converse holds too:
\par\smallskip
\textbf{Theorem.} A free-choice system is live iff it is place-live.
\par\smallskip
\textit{Proof of the non-trivial direction.} Take any reachable $M$ and target $t$ with ${}^\bullet t = \{p_1, \ldots, p_n\}$. Place-liveness yields $M \to^\ast M_1$ marking $p_1$; then $M_1 \to^\ast M_2$ marking $p_2$, and so on. The crux: the token on $p_1$ cannot vanish while reaching $M_2$: any transition $t'$ consuming it has $t' \in p_1^\bullet$, hence ${}^\bullet t' = {}^\bullet t$ by free choice, so if $t'$ was enabled then $t$ was \emph{already} enabled and we are done. Iterating, after $n$ steps all of ${}^\bullet t$ is marked and $t$ is enabled. \hfill$\square$
\par\smallskip
Liveness, a property of transitions, becomes checkable on places: the first payoff of the restriction, and the stepping stone to:
\par\smallskip
\textbf{Theorem (Commoner).} A free-choice system is live iff every proper siphon contains an initially marked trap.
\par\smallskip
Three questions remain: what is a siphon, what is a trap, and is the check feasible?
\end{notebox}

% ── slides p.22–p.33 ──────────────────────────────────────────
\subsection{Stable sets of markings}

\begin{definitionbox}{Stable set}
A set of markings $\mathbf{M}$ is \textbf{stable} if $M \in \mathbf{M} \Rightarrow [M\rangle \subseteq \mathbf{M}$: once inside, firing cannot leave. By induction it suffices to check closure under \emph{one} step: $(M \in \mathbf{M} \wedge M \xrightarrow{t} M') \Rightarrow M' \in \mathbf{M}$. The reachability set $[M_0\rangle$ is the least stable set containing $M_0$.
\end{definitionbox}

Worked candidates, all stable: $\{\mathbf{0}\}$ (the empty marking enables nothing); the set of all markings (the universe); the live markings (liveness is closed under reachability); the deadlock markings (vacuously: nothing fires); the level set $\{M \mid \mathbf{I} \cdot M = \mathbf{I} \cdot M_0\}$ of any S-invariant (the fundamental property of invariants). A non-example: $\{M \mid M(P) = 1\}$, the one-token markings: a single transition $t : p_1 \to p_2 + p_3$ jumps from one token to two.

% ── slides p.34–p.45 ──────────────────────────────────────────
\subsection{Siphons}

\begin{definitionbox}{Siphon}
$R \subseteq P$ is a \textbf{siphon} if ${}^\bullet R \subseteq R^\bullet$: every transition that feeds $R$ also consumes from $R$; \emph{proper} if non-empty. Consequence: \textbf{once empty, always empty}: the set $\{M \mid M(R) = 0\}$ is stable, since with $R$ empty no transition of ${}^\bullet R$ is enabled.
\end{definitionbox}

Worked check on the producer/consumer net, with $R = \{\textsf{prod\_busy}, \textsf{prod\_free}, \textsf{item\_buffer}\}$:
\[
  {}^\bullet R = \{\textsf{prod\_start}, \textsf{prod\_end}\}
  \subseteq
  R^\bullet = \{\textsf{prod\_start}, \textsf{prod\_end}, \textsf{cons\_start}\},
\]
so $R$ is a proper siphon. Dropping \textsf{prod\_free} breaks it: \textsf{prod\_start} still feeds \textsf{prod\_busy} but no longer consumes from the set: the three-place set was the minimal siphon around those places.

\begin{notebox}{Siphons against liveness}
\textbf{Proposition.} In a live system every proper siphon is marked at $M_0$. \textit{Proof}: pick $p \in R$ with $p^\bullet \neq \emptyset$ (a siphon that ever holds a token has one, else it could never be emptied) and $t \in p^\bullet$; liveness fires $t$ somewhere reachable, which requires $M(p) \geq 1$, so $R$ is marked there, and a siphon empty at $M_0$ would have stayed empty. \hfill$\square$ Contrapositive: an initially unmarked proper siphon is a purely structural witness of non-liveness.
\end{notebox}

% ── slides p.46–p.59 ──────────────────────────────────────────
\subsection{Traps}

\begin{definitionbox}{Trap}
$R \subseteq P$ is a \textbf{trap} if $R^\bullet \subseteq {}^\bullet R$: every transition draining $R$ also refills it; \emph{proper} if non-empty. Dual consequence: \textbf{once marked, always marked}: $\{M \mid M(R) > 0\}$ is stable.
\end{definitionbox}

Dual worked check: $R = \{\textsf{item\_buffer}, \textsf{cons\_busy}, \textsf{cons\_free}\}$ has $R^\bullet = \{\textsf{cons\_start}, \textsf{cons\_end}\} \subseteq {}^\bullet R = \{\textsf{prod\_end}, \textsf{cons\_start}, \textsf{cons\_end}\}$: a trap. Dropping \textsf{cons\_free} breaks it (\textsf{cons\_end} then drains without refilling).

\begin{notebox}{Maximal traps and Commoner, operational form}
\textbf{Lemma.} The union of traps is a trap ($(X_1 \cup X_2)^\bullet = X_1^\bullet \cup X_2^\bullet \subseteq {}^\bullet X_1 \cup {}^\bullet X_2 = {}^\bullet(X_1 \cup X_2)$). Hence every set of places contains a unique \emph{maximal} trap, and a siphon contains a marked trap iff its maximal trap is marked.
\par\smallskip
\textbf{Commoner, restated.} A free-choice system is live iff every proper siphon has its maximal trap initially marked; \emph{non-live} iff some proper siphon has an unmarked maximal trap: a single-witness formulation, the basis of the algorithmics.
\end{notebox}

% ── slides p.60–p.83 ──────────────────────────────────────────
\subsection{Complexity of liveness: NP-complete}

Non-liveness of a free-choice system is in \textsf{NP}: \emph{guess} $R \subseteq P$ (non-deterministic, polynomial); \emph{check} ${}^\bullet R \subseteq R^\bullet$ (polynomial); \emph{compute} the maximal trap $Q \subseteq R$: polynomial by the routine
\begin{lstlisting}
Q := R
while (exists p in Q, exists t in p^bullet, t notin ^bullet Q):
    Q := Q \ {p}
return Q
\end{lstlisting}
(each iteration removes a place whose consumer never refills $Q$; the residue is the largest trap inside $R$), and \emph{answer} non-live iff $M_0(Q) = 0$. A deterministic algorithm would face $2^{|P|}$ candidate subsets.

\begin{notebox}{NP-hardness: encoding CNF-SAT}
Given a CNF formula $\phi$ over variables $x_1, \ldots, x_n$ (a conjunction of clauses, each a disjunction of literals), build a free-choice net from the \emph{negation} $\neg\phi$ (a disjunction of conjunctions, by De Morgan):
\begin{itemize}
  \item one initially marked place $L_i$ per variable; two transitions $x_i, \overline{x_i}$ both consuming exactly $\{L_i\}$ (equal pre-sets: free choice preserved): firing one commits the truth value;
  \item one place per \emph{literal occurrence} in a clause of $\neg\phi$, fed by the literal's transition;
  \item one transition $C_j$ per clause of $\neg\phi$, consuming all its occurrence places: enabled iff the chosen assignment satisfies that conjunct;
  \item a place \textsf{True} fed by every $C_j$, and a transition \textsf{Back} consuming \textsf{True} and refilling every $L_i$, closing the cycle.
\end{itemize}
The net is free-choice, of size polynomial in $|\phi|$. If the chosen assignment satisfies $\phi$, it falsifies $\neg\phi$: no $C_j$ is enabled, \textsf{True} never marks, \textsf{Back} dies, the $L_i$ are never refilled: the system is \textbf{not live}. Conversely a dead run traces back to a falsifying assignment of $\neg\phi$, i.e.\ a satisfying one of $\phi$:
\[
  \phi \text{ satisfiable} \iff \text{the net is non-live}.
\]
\textbf{Verdict}: non-liveness of free-choice systems is \textsf{NP}-complete; no polynomial decision procedure for liveness exists unless $\textsf{P} = \textsf{NP}$.
\end{notebox}

% ── slides p.85–p.102 ──────────────────────────────────────────
\subsection{Liveness and boundedness: the Rank Theorem}

\begin{definitionbox}{Cluster}
For a node $x$ of $N = (P, T, F)$, the \textbf{cluster} $[x]$ is the least set with: (1) $x \in [x]$; (2) if a place is in the cluster, \emph{all its consumers} are ($p \in [x] \Rightarrow p^\bullet \subseteq [x]$); (3) if a transition is in the cluster, \emph{all its feeders} are ($t \in [x] \Rightarrow {}^\bullet t \subseteq [x]$). Forward from places, backward from transitions, to closure. Clusters partition the nodes; each cluster gathers a maximal block of intertwined conflicts. Example: in a net where $p_2^\bullet = \{t_2, t_4\}$, $p_3^\bullet = \{t_3, t_4\}$, ${}^\bullet t_3 = \{p_3, p_4\}$, the closure of $p_2$ is $[p_2] = \{p_2, p_3, p_4, t_2, t_3, t_4\}$, while an untangled pair forms its own cluster $\{p_1, t_1\}$.
\end{definitionbox}

\begin{notebox}{Rank Theorem (proof omitted)}
A free-choice system is \textbf{live and bounded} iff: (1) it has at least one place and one transition; (2) it is connected; (3) $M_0$ marks every proper siphon; (4) it has a positive S-invariant; (5) it has a positive T-invariant; (6) $\mathrm{rank}(\mathbf{N}) = |C_N| - 1$, with $C_N$ the set of clusters.
\end{notebox}

Each condition is checkable in polynomial time: (1)--(2) trivially; (4)--(5) are linear-programming feasibility; (6) is a matrix rank plus a graph closure; and (3) reduces to the dual of the trap routine: start from the initially \emph{unmarked} places $R$ and extract the maximal siphon:
\begin{lstlisting}
Q := R
while (exists p in Q, exists t in ^bullet p, t notin Q^bullet):
    Q := Q \ {p}
return Q
\end{lstlisting}
$M_0$ marks every proper siphon iff this returns $\emptyset$. Hence:

\begin{notebox}{Live and bounded is in P}
Deciding whether a free-choice system is \emph{live and bounded} takes polynomial time: in sharp contrast with plain liveness, which is \textsf{NP}-complete. Adding the boundedness requirement \emph{collapses} the complexity. The four headline results of the chapter: liveness $=$ place-liveness; non-liveness $\iff$ a proper siphon empty at some reachable marking; Commoner's siphon--trap characterisation of liveness; the Rank Theorem with its polynomial decision procedure.
\end{notebox}

% ── slides p.108–p.125 ──────────────────────────────────────────
\begin{examplebox}{Exercises}
\emph{(a) SAT via nets.} For $\phi = x_2 \wedge (x_1 \vee \overline{x_3} \vee \overline{x_4}) \wedge (x_1 \vee \overline{x_2}) \wedge (\overline{x_1} \vee x_4) \wedge (\overline{x_2} \vee \overline{x_4})$, negate, build the free-choice net of $\neg\phi$, and check liveness. The constructed net turns out \emph{live}, so $\phi$ is \textbf{not satisfiable}.
\par\smallskip
\emph{(b) Witnessing siphon.} A given free-choice system is non-live; find a proper siphon with no marked trap. Candidate $R = \{p_1, p_2, p_3, p_4, p_5, p_8\}$: the siphon condition checks out, and the maximal-trap routine strips the set step by step: $p_2, p_3$ go first (their consumers $t_2, t_3$ never refill the candidate), then $p_1$ (drained by $t_1$), then $p_8$ (drained by $t_9$), finally $p_4, p_5$, terminating at $Q = \emptyset$. The only trap inside $R$ is empty: by Commoner, non-live, with $R$ as the witness.
\par\smallskip
\emph{(c) Mirror direction.} A live free-choice system has exactly two proper siphons, $R_1 = \{p_1, \ldots, p_5, p_7, p_8\}$ and $R_2 = \{p_1, \ldots, p_5, p_6, p_8\}$. Running the routine on $R_1$ removes only $p_2$; the residue is a trap, initially marked. Symmetrically for $R_2$. Every proper siphon contains a marked trap, in agreement with Commoner and the system's liveness.
\end{examplebox}
